% Plan: Section Overview - We describe our zero-shot transfer framework, building on a pretrained DRL model. The core idea is to extract the frozen encoder and scorer, then drive them through an external multi-crane environment with a pluggable crane assignment strategy.
\section{Proposed Method: Zero-Shot Transfer Framework}
\label{sec:method}

% Plan para 1: We first describe the source DRL model we build upon (Shin et al. 2026), covering the encoder, decoder, and training. This establishes what we extract and why.
\subsection{Source DRL Model}

We adopt the publicly available pretrained policy network of Shin et al.~\cite{shin2026} as our source model. The network $\pi_\theta(a_t|s_t)$ employs a Transformer-based encoder-decoder architecture~\cite{vaswani2017} trained via REINFORCE~\cite{williams1992} with POMO~\cite{kwon2020} parallel rollouts. This section describes its architecture, as the zero-shot extraction procedure in Section~\ref{sec:extraction} directly leverages its internal representations and scoring mechanism.

\paragraph{Encoder.} The encoder processes the yard state $s_t \in \mathbb{R}^{B \times R \times T}$ through three stages. First, each stack's vertical container sequence (priority values ordered from bottom to top) is encoded via a single-layer LSTM with hidden dimension 128 and positional encoding to provide height information. The LSTM's mean-pooled output and final hidden state are concatenated and linearly projected to the embedding dimension $d = 128$. Second, the resulting stack embeddings pass through $L = 3$ multi-head self-attention layers with $H = 8$ independent heads~\cite{vaswani2017}. Each layer computes:
\begin{equation}
    \text{MultiHead}(Q,K,V) = [\text{head}_1; \dots; \text{head}_H]W_O,
\end{equation}
followed by layer normalization, a feed-forward network with hidden size 512, ReLU activation, and residual skip connection. Third, stack embeddings are aggregated at the bay level via mean pooling: $\mathbf{b}_k = \frac{1}{R} \sum_{i \in \text{bay}_k} \mathbf{h}_i$. The bay embedding is concatenated to each stack embedding and projected through an MLP: $\mathbf{h}_i' = [\mathbf{h}_i; \mathbf{b}_{k_i}] W_{\text{proj}}$. The global graph embedding is computed as $\mathbf{g} = \frac{1}{S} \sum_{i=1}^{S} \mathbf{h}_i'$, where $S = B \times R$ is the total number of stacks. Additionally, ten hand-crafted features per stack (height, remaining tiers, row and bay positions, travel distances to target, fill ratio, and layout ratio) are computed and concatenated to the learned embeddings.

\paragraph{Decoder.} The decoder computes action probabilities via cross-attention. Given the encoder outputs $\{\mathbf{h}_i'\}_{i=1}^{S}$, the global embedding $\mathbf{g}$, and the current target stack embedding $\mathbf{h}_{\text{tgt}}'$, the context vector is constructed as:
\begin{equation}
    \mathbf{c} = W_{\text{target}} \mathbf{h}_{\text{tgt}}' + W_{\text{global}} \mathbf{g}.
\end{equation}
Multi-head attention attends over node keys and values: $\text{head} = \text{MHA}(\mathbf{c}, \{W_{K1}\mathbf{h}_i'\}, \{W_V\mathbf{h}_i'\})$. The query projection $\mathbf{q} = W_Q(\text{head})$ is combined with node keys $W_{K2}\mathbf{h}_i'$ via dot-product attention, scaled by $1/\sqrt{d}$, and passed through a tanh clipping operation with parameter $C = 10$:
\begin{equation}
    u_i = C \cdot \tanh\!\Big( \frac{\mathbf{q}^\top W_{K2}\mathbf{h}_i'}{\sqrt{d}} \Big).
\end{equation}
Invalid destinations (full stacks and the target stack itself) are masked by subtracting a large constant. The final action probabilities are obtained via softmax:
\begin{equation}
    \pi_\theta(a_t = i | s_t) = \frac{\exp(u_i)}{\sum_{j \in \mathcal{V}} \exp(u_j)},
\end{equation}
where $\mathcal{V}$ is the set of feasible destination stacks.

\paragraph{Training.} The model was trained using the REINFORCE algorithm~\cite{williams1992} with POMO~\cite{kwon2020} parallel rollouts ($M = 16$). For each training instance, $M$ trajectories are generated by shifting the initial target stack across the $M$ highest-priority containers. The proposed baseline normalizes both mean and variance across rollouts: $A_m = (R_m - \mu(R)) / (\sigma(R) + \epsilon)$, where $\mu(R)$ and $\sigma(R)$ are the mean and standard deviation of rollout costs. Optimization uses Adam with learning rate $3 \times 10^{-4}$ and gradient clipping at maximum norm 1.0, for 1,000 epochs. The model achieves a verified 7.8\% mean gap against the theoretical lower bound on the full 36-instance Lee benchmark~\cite{lee2010}, and the pretrained weights are publicly available.

% Plan para 2: Our core contribution - extracting the encoder and scorer from the frozen model, enabling decoupled inference in external multi-crane environments.
\subsection{Policy Extraction and Decoupled Inference}
\label{sec:extraction}

The original model's forward pass couples the encoder-decoder policy with an internal single-crane environment loop. To enable multi-crane operation, we extract the encoder and scoring mechanism as standalone modules, forming a \texttt{ZeroShotPolicy} that exposes per-step action selection to an external environment controller.

\paragraph{Extraction.} From the pretrained parameters $\theta^*$, we extract:
\begin{enumerate}
    \item The \textbf{encoder} $\mathcal{E}$, comprising the LSTM stack encoder, multi-head self-attention layers, bay embedding MLP, and feature projection layers.
    \item The \textbf{scorer weights}: context projection matrices $W_{\text{target}}, W_{\text{global}}$; node key/value projections $W_{K1}, W_V$; query and key projections $W_Q, W_{K2}$; and the cross-attention module MHA.
    \item The \textbf{hyperparameters}: embedding dimension $d = 128$, tanh clipping $C = 10$, and cost parameters $(t_{\text{acc}}, t_{\text{bay}}, t_{\text{row}}, t_{\text{pd}})$.
\end{enumerate}

All extracted weights are frozen. No fine-tuning or gradient updates are performed, ensuring that the original model's learned representations are preserved exactly.

\paragraph{Decoupled inference.} At each decision step within a multi-crane episode, the policy computes:
\begin{align}
    \{\mathbf{h}_i'\}, \mathbf{g} &\gets \mathcal{E}(s_t) \label{eq:enc} \\
    \mathbf{c} &\gets W_{\text{target}} \mathbf{h}_{\text{tgt}}' + W_{\text{global}} \mathbf{g} \label{eq:ctx} \\
    \mathbf{q} &\gets W_Q \big( \text{MHA}([\mathbf{c}; W_{K1}\mathbf{h}'; W_V\mathbf{h}']) \big) \label{eq:query} \\
    u_i &\gets C \cdot \tanh\!\big( \mathbf{q}^\top W_{K2} \mathbf{h}_i' / \sqrt{d} \big) \label{eq:logits} \\
    d &\gets \arg\max_{i \in \mathcal{V}} u_i \label{eq:action}
\end{align}

The policy selects a destination stack $d$. This destination is then passed to a strategy module that determines which crane $c \in \{1, \dots, C\}$ executes the relocation. The complete inference procedure is specified in Algorithm~\ref{alg:zs}.

\begin{algorithm}[h!]
    \caption{Zero-shot M-CRP inference procedure. The frozen encoder and scorer select a destination stack; the strategy module assigns a crane; the multi-crane environment executes the action and returns the updated state.}
    \label{alg:zs}
    \begin{algorithmic}[1]
        \Require Frozen encoder $\mathcal{E}$, frozen scorer $\mathcal{D}$, multi-crane environment $\mathcal{E}_{\text{MC}}$, assignment strategy $S$, number of cranes $C$
        \Ensure Total working time $J$
        \State $s \gets \mathcal{E}_{\text{MC}}.\text{reset}()$
        \State $J \gets 0$
        \While{$\neg \mathcal{E}_{\text{MC}}.\text{terminated}$}
            \State $\{\mathbf{h}_i'\}, \mathbf{g} \gets \mathcal{E}(s)$ \Comment{Encode yard state}
            \State $\{u_i\} \gets \mathcal{D}(\{\mathbf{h}_i'\}, \mathbf{g}, s.\text{target})$ \Comment{Compute relocation scores}
            \State $d \gets \arg\max_i u_i$ \Comment{Select destination stack}
            \State $c \gets S.\text{assign}(\mathcal{E}_{\text{MC}}, d)$ \Comment{Select executing crane}
            \State $(\Delta J, s) \gets \mathcal{E}_{\text{MC}}.\text{step}(d, c)$ \Comment{Execute relocation}
            \State $J \gets J + \Delta J$
        \EndWhile
    \end{algorithmic}
\end{algorithm}

% Plan para 3: The four crane assignment strategies span the design space from simple to complex. S1 and S3 are spatially unaware; S2 and S4 exploit spatial information.
\subsection{Crane Assignment Strategies}
\label{sec:strategies}

We design four assignment strategies (Table~\ref{tab:strategies}) that span the design space from $O(1)$ cyclic assignment to $O(C^2)$ one-step lookahead. Strategies S1 and S3 are \emph{spatially unaware}--they assign tasks to cranes without reference to crane positions or task locations. Strategies S2 and S4 are \emph{spatially aware}--they exploit crane position information to minimize interference and travel costs. This separation allows us to isolate the effect of spatial coordination on zero-shot transfer quality.

\begin{table}[h!]
    \centering
    \caption{The four crane assignment strategies for zero-shot M-CRP transfer. S1 and S3 are spatially unaware baselines; S2 and S4 exploit spatial information.}
    \label{tab:strategies}
    \begin{tabularx}{\textwidth}{lXcc}
        \toprule
        Strategy & Mechanism & Spatial? & Complexity \\
        \midrule
        S1: RoundRobin & Cyclic: crane $(t \bmod C)$ receives task $t$ & No & $O(1)$ \\
        S2: ZoneSplit & Static bay zones: each crane serves a fixed bay range & Yes & $O(B)$ \\
        S3: LoadBalance & Min-count: assign to crane with fewest completed tasks & No & $O(C)$ \\
        S4: GreedyOptimal & 1-step lookahead: minimize projected travel + interference & Yes & $O(C^2)$ \\
        \bottomrule
    \end{tabularx}
\end{table}

\paragraph{S1: RoundRobin.} Tasks are assigned in a fixed cyclic order, $\text{crane}(t) = (t \bmod C)$. This strategy ignores both the spatial location of the destination stack and the current positions of all cranes. It serves as a lower bound on coordination quality.

\paragraph{S2: ZoneSplit.} The yard is divided into $C$ contiguous bay zones. Crane $c$ is assigned tasks whose source stack (the target stack requiring relocation) falls within its designated zone. Zone boundaries are computed as:
\begin{equation}
    z_c^{\text{start}} = \left\lfloor \frac{(c-1) \cdot B}{C} \right\rfloor + 1, \quad
    z_c^{\text{end}} = \left\lfloor \frac{c \cdot B}{C} \right\rfloor.
\end{equation}
This strategy enforces spatial separation by construction: since each crane operates in a distinct set of bays, constraints A6 (spatial non-overlap) and A7 (non-crossing) are automatically satisfied. ZoneSplit requires $O(B)$ computation to determine zone membership and is compatible with existing terminal zoning practices.

\paragraph{S3: LoadBalance.} Tasks are assigned to the crane with the fewest completed tasks at each decision step: $c = \arg\min_c \text{task\_count}[c]$. This strategy balances the total number of relocations across cranes without considering the spatial cost of those relocations. It isolates the effect of load balancing from spatial coordination.

\paragraph{S4: GreedyOptimal.} For each crane $c$, we compute the projected cost of executing the relocation, including travel from the crane's current position and an interference penalty for destination bays occupied by other cranes:
\begin{equation}
    \text{cost}(c) = \text{travel}(\text{pos}_c, \text{bay}_d) + \sum_{c' \neq c} \mathbb{I}[\text{bay}_{c'} = \text{bay}_d] \cdot t_{\text{acc}}.
\end{equation}
The crane minimizing this projected cost is selected. This represents a single-step lookahead that balances travel efficiency with interference avoidance, at $O(C^2)$ cost due to the pairwise interference checks.

% Plan para 4: The multi-crane environment extends the single-crane CRP environment with independent crane tracking and interference validation.
\subsection{Multi-Crane Environment}
\label{sec:mcenv}

The \texttt{MCEnv} class wraps the single-crane CRP environment~\cite{shin2026} to support $C$ simultaneously tracked cranes. Key design elements include:

\begin{enumerate}
    \item \textbf{Independent position tracking:} Each crane maintains its own current bay and row position, initialized from a user-specified or automatically distributed start configuration.
    \item \textbf{Crane busy state:} Each crane has a boolean busy flag preventing simultaneous task assignment, ensuring that exactly one task executes per decision step across all cranes.
    \item \textbf{Interference validation:} Before execution, the destination bay is checked against the current positions of all cranes. If constraint A6 is violated (destination bay occupied by another crane), interference resolution reassigns the task to an alternative crane.
    \item \textbf{Cost delegation:} The base environment's \texttt{curr\_bay} and \texttt{curr\_row} attributes are set to the executing crane's position before cost calculation, ensuring that travel costs reflect the correct crane origin. After execution, the crane's position is updated to the destination bay and row.
\end{enumerate}

The total per-step computational cost comprises the DRL encoder and scorer ($\approx$5 ms on CPU) plus the strategy assignment (negligible, $O(1)$ to $O(C^2)$). The combined cost ($\approx$5 ms) is suitable for real-time control applications.
